\subsection{PRESENTATIONS OF A LIE ALGEBRA}

Let \(\mathfrak{g}\) be a Lie algebra and \(\pmb{a} = (a_i)_{i\in I}\) a family of elements of \(\mathfrak{g}\) . Let \(f_{a}\) be the homomorphism of \(\mathrm{L(I)}\) into \(\mathfrak{g}\) mapping each \(i\in I\) to \(a_{i}\) . The image of \(f_{a}\) is the subalgebra of \(\mathfrak{g}\) generated by \(\pmb{a}\) ; the elements of the kernel of \(f_{a}\) are called the relators of the family \(\pmb{a}\) . The family \(\pmb{a}\) is called generating (resp. free, basic) if \(f_{a}\) is surjective (resp. injective, bijective).

Let \(\mathfrak{g}\) be a Lie algebra. A presentation of \(\mathfrak{g}\) is an ordered pair \((\boldsymbol{a},\boldsymbol{r})\) consisting of a generating family \(\boldsymbol{a} = (a_{i})_{i\in \mathbb{I}}\) and a family \(\boldsymbol{r} = (r_j)_{j\in \mathbb{J}}\) of relators of \(\boldsymbol{a}\) generating the ideal of L(I) the kernel of \(f_{a}\) . We also say that \(\mathfrak{g}\) is presented by the family \(\boldsymbol{a}\) related by the relators \(r_j\) ( \(j\in J\) ).

Let I be a set and \(\pmb{r} = (r_j)_{j \in J}\) be a family of elements of the free Lie algebra L(I); let \(a_r\) be the ideal of L(I) generated by \(\pmb{r}\) . The quotient algebra L(I, \(\pmb{r}\) ) = L(I)/ \(a_r\) is called the Lie algebra defined by I and the family of relators ( \(r_j)_{j \in J}\) ; we also say that L(I, \(\pmb{r}\) ) is defined by the presentation (I, \(\pmb{r}\) ), or also by (I; ( \(r_j = 0\) ) \(_{j \in J}\) ). When the family \(\pmb{r}\) is empty, L(I, \(\pmb{r}\) ) = L(I).

Let \(I\) and \(r\) be as above; let \(\xi_i\) denote the image of \(i\) in \(L(I, r)\) . The generating family \(\xi = (\xi_i)_{i \in I}\) and the family of relators \(r\) constitute a presentation of \(L(I, r)\) . Conversely, if \(g\) is a Lie algebra and \((a, r)\) , where \(a = (a_i)_{i \in I}\) , is a presentation of \(g\) , there exists a unique isomorphism \(u: L(I, r) \to g\) such that \(u(\xi_i) = a_i\) for all \(i \in I\) .

